\section{AI 产品像挖矿：Magic Moment 与保鲜窗口}

本章解释“挖矿”比喻。第一个做出让用户惊叹体验的产品，会获得巨大营销红利，即使 token 消耗很高也值得。早期搜索问答产品、Cursor、Manus 都属于抓住了某个 Magic Moment 的例子。但窗口期正在缩短：从两年、一年，到几个月。

\begin{figure}[H]
\centering
\includegraphics[width=0.96\textwidth]{ai-product-mining-window.png}
\caption{AI 产品像挖矿：Magic Moment 的保鲜窗口正在缩短。自制概念图，依据 00:38:52--00:44:21 对谈内容整理。}
\end{figure}

\begin{knowledgebox}{读图：挖矿不是长期护城河}
挖矿阶段抢的是第一次惊叹和品牌红利；窗口缩短后，模型公司、竞品和开源实现都会迅速跟进。真正护城河要从 Magic Moment 过渡到工作流、数据、入口和品牌。
\end{knowledgebox}

\subsection{产品公司能赢过模型公司吗}

挖矿窗口的关键问题是：产品公司能不能在模型公司下场前形成壁垒。产品公司有速度、专注和体验优势；模型公司有模型、分发、算力和全家桶优势。产品公司要赢，必须在窗口期内形成强用户心智和数据闭环，否则会被模型公司功能化吸收。

\begin{importantbox}{窗口期公式}
\[
\text{Startup chance} \approx \text{Magic Moment} \times \text{Speed} \times \text{Moat after copy}
\]
其中，Magic Moment 是首次惊叹，Speed 是抢窗口速度，Moat after copy 是被复制后仍然保留的壁垒。
\end{importantbox}

\begin{knowledgebox}{窗口缩短的原因}
窗口期缩短有三个原因。第一，基础模型能力更新更快，旧体验很快变成默认能力。第二，头部模型公司有更强分发，能迅速复制成功模式。第三，开源和工具链让工程实现更快，产品领先不能只靠 demo。
\end{knowledgebox}

\begin{knowledgebox}{创业者策略表}
\begin{tabularx}{\textwidth}{p{0.22\textwidth}p{0.34\textwidth}X}
\toprule
策略 & 适用场景 & 风险 \\
\midrule
抢 Magic Moment & 新能力刚出现，用户还没形成默认入口 & 窗口短，容易被模型公司复制。 \\
垂直深扎 & 行业工作流复杂，数据和客户关系难迁移 & 增长慢，销售和交付重。 \\
做工具链 & 围绕头部模型做开发者/企业工具 & 平台 API 变化会影响产品。 \\
做入口 & 抢默认使用场景和品牌心智 & 与 ChatGPT/Gemini 正面竞争。 \\
\bottomrule
\end{tabularx}
\end{knowledgebox}

\subsection{本章小结}

AI 产品像挖矿，是因为新能力出现时会有短暂高收益窗口。创业者要做的不是只挖第一铲，而是在窗口关闭前形成长期壁垒。

